\def\subsectionname{Convergence}
\subsection{\subsectionname}

Une suite bornée n'est pas forcément convergente.
% un commentaire seul ne coupe pas le paragraphe
Mais elle a une sous-suite qui l'est : c'est le résultat central.

\begin{theorem}[title={Bolzano-Weierstrass}, label=thm:bw]
  Toute suite réelle bornée a une sous-suite convergente.
\end{theorem}%
\footnotetext{Une note de l'énoncé.}%
\begin{proof}
  Par dichotomie.
\end{proof}

On en tire la compacité des segments.

\begin{remark}
  Une remarque, avec un exemple dedans :
  \begin{example}
    $u_n = (-1)^n$.
  \end{example}
\end{remark}
\begin{itemize}
  \item une liste fait partie du texte ;
\end{itemize}
Voici l'exercice.

\begin{exercise}
  Montrer que $(-1)^n$ a deux valeurs d'adhérence.
\end{exercise}

\section{Suite}
Le cadre de la section.

\begin{assumption}[label=hyp:h]
  $u$ est bornée.
\end{assumption}
